\documentclass[11pt,twocolumn]{article}
\usepackage[top=1in,bottom=1in,left=0.85in,right=0.85in]{geometry}
\usepackage{amsmath,amssymb,amsthm}
\usepackage{natbib}
\usepackage{graphicx}
\usepackage{hyperref}
\usepackage{booktabs}


\theoremstyle{plain}
\newtheorem{theorem}{Theorem}[section]
\newtheorem{lemma}[theorem]{Lemma}
\newtheorem{corollary}[theorem]{Corollary}
\theoremstyle{definition}
\newtheorem{definition}[theorem]{Definition}
\newtheorem{assumption}[theorem]{Assumption}

\title{Adaptive-Regret Bed Allocation: \\
A CUSUM-Triggered Robust Optimization Framework \\
for Non-Stationary Healthcare Environments}

\author{Darshan Sathish Kumar \\
UTSA Honors College, Department of Mathematics and Computer Science\\
\texttt{darshan.sathishkumar@my.utsa.edu}}

\date{\today}

\begin{document}

\maketitle

\begin{abstract}
We consider the problem of allocating emergency department (ED) beds to arriving patients with varying acuity levels under non-stationary demand. Traditional robust optimization assumes a fixed uncertainty set over the planning horizon, but ED arrival rates exhibit pronounced non-stationarity: diurnal cycles, weekly seasonality, and abrupt surge events such as mass casualty incidents or infectious disease outbreaks. A static robust policy tuned for normal conditions is dangerously under-protective during surges, while one tuned for surges is wastefully conservative during normal operation. We propose an adaptive framework that couples a Poisson Cumulative Sum (CUSUM) sequential change detector with an online robust allocation policy. When the CUSUM statistic crosses a calibrated threshold, the policy re-centers its box uncertainty set via sliding-window maximum likelihood estimation and expands its robustness budget. The detection delay is bounded by Lorden's inequality, and the post-detection regret follows from standard online convex optimization bounds. We evaluate the framework on a SimPy discrete-event simulation of a 20-bed ED with NHAMCS-calibrated parameters across four scenarios --- baseline, step-shift surge, linear drift, and acuity composition shift --- using 30 stochastic replications each. The adaptive policy achieves 17--29\% cumulative regret reduction over a static robust baseline (all $p < 0.001$, paired $t$-test). Detection delays are well within Lorden's theoretical bound: 1.1~hours for step shifts, 4.8~hours for gradual drifts, with zero false alarms in normal operation. A MIMIC-IV validation pipeline demonstrates 21\% regret reduction under real-data-calibrated parameters. Our results establish that coupling sequential change detection with robust optimization provides a principled, computationally tractable approach to resource allocation under non-stationary uncertainty.
\end{abstract}

\section{Introduction}

Emergency department crowding is a persistent public health crisis with demonstrable consequences for patient outcomes. Studies consistently document that crowding correlates with higher mortality rates, prolonged inpatient stays once admitted, increased ambulance diversion, and higher rates of patients leaving without being seen \cite{sridhar2013patient,golshani2019bed}. At its operational core, ED crowding reflects a resource allocation problem: an ED has a fixed number of beds, but patient demand --- both in volume and acuity composition --- exhibits pronounced non-stationarity. Arrivals follow diurnal patterns (peaking in late morning, troughing overnight), weekly seasonality, and are punctuated by shock events such as mass casualty incidents, seasonal influenza outbreaks, or severe weather.

Robust optimization has emerged as a leading framework for decision-making under uncertainty in healthcare operations \cite{bertsimas2004price,bental2009robust,chan2016robust}. By optimizing against worst-case realizations within a prescribed uncertainty set, robust methods provide guarantees that nominal approaches lack. However, standard robust optimization assumes the uncertainty set is fixed over the planning horizon. This assumption is reasonable for stationary processes, but ED arrivals are not stationary. A static box uncertainty set centered at the average arrival rate under-protects during a surge (the true rate lies outside the set) and over-conserves during normal operation (the set is unnecessarily wide). The practical consequence is that a static robust policy incurs regret --- excess cost over a hindsight-optimal policy --- that grows linearly with the post-shift horizon.

This paper addresses the fundamental limitation by answering the following question: \emph{can a robust allocation policy detect when the underlying distribution has shifted and adapt its uncertainty set accordingly, and if so, what regret improvement can be guaranteed?}

\paragraph{Related work.} Our work bridges three distinct literatures. First, \citet{bertsimas2004price} introduced $\Gamma$-robustness, which interpolates between nominal and worst-case optimization via a budget-of-uncertainty parameter; this provides the robust optimization foundation. \citet{chan2016robust} applied robust optimization to ED staffing, and \citet{nohadani2017robust} extended the framework to time-dependent uncertainty sets in radiation therapy, demonstrating that uncertainty sets can be partitioned across temporal segments. Our work automates this segmentation using sequential change detection.

Second, sequential change-point detection dates to \citet{page1954continuous}, who introduced the CUSUM for quality control in manufacturing. \citet{lorden1971procedures,moustakides1986optimal} established CUSUM's optimality: among all stopping rules with a given in-control false-alarm rate, CUSUM minimizes the worst-case detection delay. \citet{lai1995sequential} provides a comprehensive survey. These theoretical guarantees make CUSUM the natural detector for our setting.

Third, the online convex optimization literature provides regret bounds for sequential decision-making. \citet{zinkevich2003online} proved $\mathcal{O}(\sqrt{T})$ regret for online gradient descent, and \citet{hazan2007logarithmic} established $\mathcal{O}(\log T)$ regret for exp-concave losses via the Online Newton Step. Our post-detection rate estimation falls into this framework. The closest work to ours is \citet{golrezaei2026online}, who study online robust resource allocation with adaptive regret in a stationary setting; our contribution is extending their framework to non-stationary environments via change detection.

\paragraph{Contribution.} We make four contributions:
\begin{enumerate}
\item A novel framework combining CUSUM change detection with online robust bed allocation, where detection triggers re-centering and expansion of the uncertainty set.
\item A theoretical bound on detection delay via Lorden's inequality, establishing that the expected delay is $\mathcal{O}(h / D_{\text{KL}})$ for threshold $h$ and KL divergence $D_{\text{KL}}$.
\item Empirical demonstration of 17--29\% regret reduction across four non-stationary scenarios using 30 stochastic replications, with statistical significance at $p < 0.001$.
\item A MILP-based hindsight oracle using HiGHS that provides tighter lower bounds than the greedy heuristic used in prior work.
\end{enumerate}

All code, experiments, and documentation are publicly available at \texttt{github.com/dshan12/arba}.

\section{Problem Formulation}

We model the ED as a discrete-time Markov decision process (MDP) \cite{puterman1994markov} with finite horizon $T$. Decision epochs correspond to patient arrival events.

\subsection{State Space}

At time $t$, the state $s_t \in \mathcal{S}$ captures all information relevant to the allocation decision:
\begin{equation}
s_t = (B_{\text{occ}}, q_{\text{crit}}, q_{\text{urg}}, q_{\text{rout}}, \hat{\lambda}_t, S_t),
\label{eq:state}
\end{equation}
where $B_{\text{occ}} \in \{0, 1, \dots, B_{\max}\}$ is the number of occupied beds, $q_i \in \mathbb{Z}_{\ge 0}$ are queue lengths for each acuity level $i \in \{\text{crit}, \text{urg}, \text{rout}\}$, $\hat{\lambda}_t$ is the current maximum-likelihood estimate of the arrival rate, and $S_t$ is the CUSUM statistic (Section~\ref{sec:cusum}). The state space is countably infinite due to unbounded queues, but in practice the finite horizon and bounded arrival rate keep it manageable.

\subsection{Action Space}

At each decision epoch, the policy chooses how many patients of each acuity to admit:
\begin{equation}
a_t = (a_{\text{crit}}, a_{\text{urg}}, a_{\text{rout}}) \in \mathcal{A}(s_t),
\end{equation}
where $\sum_i a_i \le B_{\max} - B_{\text{occ}}$ and $a_i \le q_i$. The action space size is $\mathcal{O}((B_{\max})^3)$ and is enumerated explicitly.

\subsection{Patient Acuity and Costs}

We use the Emergency Severity Index (ESI) aggregated to three levels, following standard practice in healthcare operations research:
\begin{itemize}
\item \textbf{Critical} (ESI 1-2): weight $w_{\text{crit}} = 100$, service time $\sim \operatorname{Lognormal}(\mu=3\text{h}, \sigma=1\text{h})$
\item \textbf{Urgent} (ESI 3): weight $w_{\text{urg}} = 10$, service time $\sim \operatorname{Lognormal}(\mu=2\text{h}, \sigma=0.75\text{h})$
\item \textbf{Routine} (ESI 4-5): weight $w_{\text{rout}} = 1$, service time $\sim \operatorname{Lognormal}(\mu=1\text{h}, \sigma=0.5\text{h})$
\end{itemize}

The weights reflect clinical priority: a 1-hour delay for a critical patient is treated as equivalent to a 100-hour delay for a routine patient. Service times follow log-normal distributions calibrated to published NHAMCS length-of-stay statistics \cite{johnson2023mimic}.

\subsection{Arrival Process}

Arrivals follow a non-homogeneous Poisson process (NHPP) with diurnal periodicity and optional surge:
\begin{equation}
\lambda(t) = \lambda_{\text{base}} \left(1 + A \sin\left(\frac{2\pi t}{24} + \phi\right)\right) \cdot f_{\text{surge}}(t),
\label{eq:nhpp}
\end{equation}
with $A = 0.3$, $\phi = \pi/2$ (peak at midnight, trough at noon). The surge function $f_{\text{surge}}(t)$ takes different forms per scenario (Section~\ref{sec:scenarios}). The expected number of arrivals in an interval of length $\Delta$ is $\int_t^{t+\Delta} \lambda(\tau)\,d\tau$.

\subsection{Cost and Regret}

The instantaneous cost incurred by action $a$ in state $s$ is:
\begin{equation}
c(s, a) = \sum_{i \in \{\text{crit},\text{urg},\text{rout}\}} w_i \cdot \mathbb{E}[\text{Wait}_i \mid s, a] + P_{\text{reject}},
\label{eq:cost}
\end{equation}
where $P_{\text{reject}} = w_i \cdot T$ for patients of acuity $i$ who are never served (they wait the entire horizon). The total cost of a policy $\pi$ over horizon $T$ is $C^{\pi}(T) = \sum_{t=1}^{T} c(s_t, \pi(s_t))$.

Let $V^*_T$ be the cost of a hindsight-optimal policy that knows all future arrivals, service times, and acuities in advance. The regret of policy $\pi$ is:
\begin{equation}
R^{\pi}(T) = C^{\pi}(T) - V^*_T.
\label{eq:regret}
\end{equation}
Regret is non-negative by definition; a smaller regret indicates better performance. The hindsight oracle is described in Section~\ref{sec:oracle}.

\section{Change Detection via CUSUM}
\label{sec:cusum}

The effectiveness of an adaptive policy depends on its ability to detect distribution shifts quickly and reliably. We use the Cumulative Sum (CUSUM) procedure, which is optimal in the sense of minimizing worst-case detection delay for a given false-alarm rate \cite{lorden1971procedures,moustakides1986optimal}.

\subsection{Poisson Log-Likelihood Ratio CUSUM}

For Poisson-distributed observations $X_t \sim \text{Poisson}(\lambda)$, testing the null hypothesis $H_0: \lambda = \lambda_0$ against $H_1: \lambda = \lambda_1$ (with $\lambda_1 > \lambda_0$), the log-likelihood ratio (LLR) for a single observation $x$ is:
\begin{equation}
\ell(x) = \log\frac{f_{\lambda_1}(x)}{f_{\lambda_0}(x)} = x \log\frac{\lambda_1}{\lambda_0} - (\lambda_1 - \lambda_0).
\label{eq:llr}
\end{equation}

The CUSUM statistic accumulates positive evidence for a change:
\begin{equation}
S_0 = 0, \qquad S_t = \max\left(0, S_{t-1} + \ell(X_t)\right).
\label{eq:cusum}
\end{equation}

A change is declared at the stopping time $\tau = \inf\{t \ge 1 : S_t \ge h\}$, where $h > 0$ is the detection threshold. The $\max(0, \cdot)$ operation resets the statistic to zero when cumulative evidence becomes negative, preventing negative drift under $H_0$.

\subsection{Reference Value and KL Divergence}

The increment $\ell(X_t)$ has expectation:
\begin{align}
\mathbb{E}_0[\ell(X_t)] &= -D_{\text{KL}}(\lambda_0 \parallel \lambda_1), \\
\mathbb{E}_1[\ell(X_t)] &= +D_{\text{KL}}(\lambda_1 \parallel \lambda_0),
\end{align}
where $D_{\text{KL}}(\lambda_1 \parallel \lambda_0) = \lambda_1 \log(\lambda_1/\lambda_0) - (\lambda_1 - \lambda_0)$ is the Kullback-Leibler divergence for Poisson distributions. Under $H_0$, the statistic drifts slowly downward and is frequently reset to zero; under $H_1$, it drifts upward at rate $D_{\text{KL}}$ per observation, eventually crossing the threshold.

The reference value $k$ --- the observation value at which the LLR is zero --- is:
\begin{equation}
k = \frac{\lambda_1 - \lambda_0}{\log(\lambda_1/\lambda_0)},
\label{eq:refval}
\end{equation}
which is a weighted average of $\lambda_0$ and $\lambda_1$ in log-space. Observations above $k$ increase $S_t$; observations below $k$ decrease it.

\subsection{Lorden's Detection Delay Bound}
\label{sec:lorden}

\begin{theorem}[Lorden's Inequality]
Let $\tau$ be the CUSUM stopping time with threshold $h$ for detecting a shift from $\lambda_0$ to $\lambda_1$. The expected detection delay under the post-change distribution satisfies:
\begin{equation}
\mathbb{E}_1[\tau] \le \frac{h + D_{\text{KL}}(\lambda_1 \parallel \lambda_0)}{D_{\text{KL}}(\lambda_1 \parallel \lambda_0)} = 1 + \frac{h}{D_{\text{KL}}(\lambda_1 \parallel \lambda_0)}.
\label{eq:lorden}
\end{equation}
\end{theorem}

\begin{proof}[Proof sketch]
Let $Y_t = \ell(X_t)$ be the LLR increments. Under $\lambda_1$, $\{Y_t\}$ are i.i.d.\ with mean $D_{\text{KL}} > 0$. At the stopping time $\tau$, we have $S_\tau \ge h$ and $S_\tau \le h + \max(0, Y_\tau)$. Wald's identity gives $\mathbb{E}_1[S_\tau] = \mathbb{E}_1[\tau] \cdot D_{\text{KL}}$, and the bound follows by combining these facts.
\end{proof}

For large $h$, the bound simplifies to $\mathbb{E}_1[\tau] \approx h / D_{\text{KL}}$, showing that detection delay grows linearly in the threshold and inversely in the KL divergence. This is the fundamental tradeoff: a higher threshold reduces false alarms but increases detection delay.

\subsection{Threshold Calibration}

The average run length in control, $\text{ARL}_0 = \mathbb{E}_\infty[\tau]$, is the expected time until a false alarm when no shift has occurred. Siegmund's approximation \cite{siegmund1985sequential} gives:
\begin{equation}
\text{ARL}_0 \approx \frac{e^{h} - h - 1}{D_{\text{KL}}(\lambda_0 \parallel \lambda_1)}.
\label{eq:arl0}
\end{equation}

We calibrate $h$ via Monte Carlo simulation. For a target $\text{ARL}_0$ (e.g., 168 hours, representing one week of operation between false alarms), we perform binary search over $h \in [1, 50]$ with 2,000 simulated Poisson streams of length 500 each. This yields the threshold used by the adaptive policy.

\section{Robust Optimization Framework}

\subsection{The Bertsimas-Sim $\Gamma$-Robustness Model}

Following \citet{bertsimas2004price}, we model arrival rate uncertainty via a box uncertainty set centered at the estimated rate $\hat{\lambda}$:
\begin{equation}
\mathcal{U}(\hat{\lambda}, \Gamma) = \{\lambda : |\lambda - \hat{\lambda}| \le \Gamma \cdot \hat{\lambda} \},
\label{eq:uncertainty}
\end{equation}
where $\Gamma \in [0, 1)$ is the robustness budget controlling the degree of conservatism. A larger $\Gamma$ produces a wider set, hedging against more severe deviations at the cost of reserving more beds for high-acuity patients.

The robust admission decision solves:
\begin{equation}
\min_{a \in \mathcal{A}(s)} \max_{\lambda \in \mathcal{U}(\hat{\lambda}, \Gamma)} \mathbb{E}_\lambda[c(s, a) \mid s],
\label{eq:robust-lp}
\end{equation}
where the inner maximization evaluates worst-case cost over the uncertainty set, and the outer minimization chooses the admission action. For our linear cost structure with priority-ordered acuity weights, the optimal solution to \eqref{eq:robust-lp} admits a closed form.

\subsection{RobustPolicy: Closed-Form Allocation}

The optimal solution to \eqref{eq:robust-lp} for our model reduces to a bed-reservation heuristic. Let $B_{\text{avail}} = B_{\max} - B_{\text{occ}}$ be the number of available beds. The robust policy reserves a fraction $\Gamma$ of available beds for critical patients:
\begin{align}
a_{\text{crit}} &= \min(q_{\text{crit}}, \lceil \Gamma \cdot B_{\text{avail}} \rceil), \\
a_{\text{urg}} &= \min(q_{\text{urg}}, B_{\text{avail}} - a_{\text{crit}} - a_{\text{rout}}), \\
a_{\text{rout}} &= \min(q_{\text{rout}}, B_{\text{avail}} - a_{\text{crit}} - a_{\text{urg}}),
\label{eq:robust-policy}
\end{align}
with remaining beds allocated to urgent then routine patients. This reservation ensures capacity is available for high-acuity arrivals even when the true arrival rate is at the upper bound of $\mathcal{U}$. For $\Gamma = 0.2$, the policy reserves 20\% of available beds for critical patients.

\subsection{AdaptivePolicy: CUSUM-Triggered Switching}

The adaptive policy augments the robust policy with CUSUM-based model switching. The algorithm is as follows:

\medskip
\noindent\fbox{\parbox{\dimexpr\columnwidth-2\fboxsep-2\fboxrule\relax}{
\small
\textbf{Algorithm: Adaptive-Regret Bed Allocation}
\vspace{2mm}

\textbf{Input:} $\lambda_0$, $\lambda_1 = \lambda_0 + 2\delta$, threshold $h$, window $w$, $\Gamma = 0.2$ \\
Initialize CUSUM($\lambda_0$, $\lambda_1$, $h$); \quad
Policy $\leftarrow$ RobustPolicy($\lambda_0$, $\Gamma$); \quad
adaptive\_mode $\leftarrow$ False

\vspace{1mm}
For each arrival epoch $t = 1, 2, \dots$:
\begin{enumerate}
  \item Observe arrival $X_t$
  \item $S_t \leftarrow \max(0, S_{t-1} + \ell(X_t))$
  \item If $S_t \ge h$ and not adaptive\_mode:
        \begin{enumerate}
          \item adaptive\_mode $\leftarrow$ True
          \item $\hat{\lambda}_{\text{new}} \leftarrow \frac{1}{w} \sum_{i=t-w}^{t} X_i$
          \item $\Gamma_{\text{eff}} \leftarrow \min(0.5, \Gamma \cdot \hat{\lambda}_{\text{new}} / \lambda_0)$
          \item CUSUM.reset($\lambda_0 = \hat{\lambda}_{\text{new}}$)
        \end{enumerate}
  \item If adaptive\_mode:
        \begin{enumerate}
          \item If expansion timer expired:
                \begin{enumerate}
                  \item $\hat{\lambda}_{\text{current}} \leftarrow 0.7\hat{\lambda}_{\text{current}} + 0.3 \cdot \frac{1}{w} \sum_{i=t-w}^{t} X_i$
                  \item If $|\hat{\lambda}_{\text{current}} - \lambda_0| < \epsilon$:
                        adaptive\_mode $\leftarrow$ False; CUSUM.reset($\lambda_0$)
                \end{enumerate}
        \end{enumerate}
  \item $a_t \leftarrow \text{Policy.get\_action}(s_t)$
\end{enumerate}
}}
\medskip

When the CUSUM triggers, the policy:
\begin{enumerate}
\item Estimates the new arrival rate $\hat{\lambda}_{\text{new}}$ via sliding-window MLE over the last $w = 6$ observations.
\item Expands the effective robustness budget: $\Gamma_{\text{eff}} = \min(0.5, \Gamma \cdot \hat{\lambda}_{\text{new}} / \lambda_0)$. This is the key adaptive mechanism: if the surge is 2x baseline, reserve up to 40\% of beds; if 3x, reserve up to 50\%.
\item Re-centers the CUSUM to detect further shifts from the new baseline.
\end{enumerate}

While in adaptive mode, the rate estimate is updated via exponential smoothing: $\hat{\lambda}_t = 0.7\hat{\lambda}_{t-1} + 0.3\bar{X}_{\text{window}}$. When the estimate returns to within $\epsilon = 0.5$ of the original baseline, the policy reverts to standard robust mode.

\section{Theoretical Analysis}
\label{sec:theory}

We analyze the regret of the adaptive policy by decomposing it into three phases: pre-detection, detection, and post-detection.

\subsection{Regret Decomposition}

Let $\tau^*$ be the true change point and $\tau$ the CUSUM detection time. The total expected regret decomposes as:
\begin{align}
\mathbb{E}[R_{\text{adaptive}}(T)] &= \mathbb{E}[R_{\text{pre}}(\tau^*)] \notag \\
&\quad + \mathbb{E}[R_{\text{detect}}(\tau - \tau^*)] \notag \\
&\quad + \mathbb{E}[R_{\text{post}}(T - \tau)].
\label{eq:decomp}
\end{align}

\textbf{Pre-detection regret} ($t < \tau^*$): Before the change, the adaptive policy coincides with the robust policy. The per-step regret is bounded by the maximum possible cost per patient times the number of beds:
\begin{equation}
\mathbb{E}[R_{\text{pre}}(\tau^*)] \le B_{\max} \cdot \max_i w_i \cdot \tau^* = 100 B_{\max} \tau^*.
\label{eq:pre-bound}
\end{equation}
This is conservative but unavoidable without knowledge of future arrivals.

\textbf{Detection regret} ($\tau^* \le t < \tau$): During the detection window, the policy operates with a stale uncertainty set centered at $\lambda_0$ while the true rate is $\lambda_1$. The regret grows linearly in the detection delay. By Lorden's inequality (Theorem~1):
\begin{equation}
\mathbb{E}[R_{\text{detect}}(\tau - \tau^*)] \le 100 B_{\max}
\Bigl(1 + \tfrac{h}{D_{\text{KL}}(\lambda_1 \parallel \lambda_0)}\Bigr).
\label{eq:detect-bound}
\end{equation}

\textbf{Post-detection regret} ($t \ge \tau$): After detection, the policy re-centers its uncertainty set around $\hat{\lambda}_{\text{new}}$ and expands $\Gamma$. The post-detection regret depends on the accuracy of the rate estimate. With sliding-window MLE over $w$ observations:
\begin{equation}
\mathbb{E}[(\hat{\lambda} - \lambda_1)^2] = \frac{\lambda_1}{w}.
\label{eq:mse}
\end{equation}
By Markov's inequality, $|\hat{\lambda} - \lambda_1| \le \delta$ with high probability when $w \ge \lambda_1 / \delta^2$. Within this regime, the robust policy with $\mathcal{U}(\hat{\lambda}, \Gamma_{\text{eff}})$ dominates the static robust policy with $\mathcal{U}(\lambda_0, \Gamma)$.

\subsection{Total Regret Bound}

\begin{theorem}[Total Adaptive Regret]
Let $\tau^*$ be the change point, $\lambda_1 > \lambda_0$ the post-change rate, and $h$ the CUSUM threshold. The expected regret of the adaptive policy over horizon $T$ is bounded by:
\begin{align}
\mathbb{E}[R_{\text{adaptive}}(T)] &\le 100 B_{\max} \tau^* \notag \\
&\quad + 100 B_{\max} \left(1 + \frac{h}{D_{\text{KL}}}\right) \notag \\
&\quad + (T - \tau^*) \cdot \mathbb{E}[r_{\text{post}}],
\label{eq:total-bound}
\end{align}
where $\mathbb{E}[r_{\text{post}}] \le \Gamma \cdot |\lambda_1 - \lambda_0| \cdot B_{\max}$ is the per-step regret after convergence, which is strictly smaller than the static robust policy's per-step regret because the uncertainty set is centered at $\lambda_1$ rather than $\lambda_0$.
\end{theorem}

The static robust policy, by contrast, never adjusts its uncertainty set. Its expected regret is:
\begin{equation}
\mathbb{E}[R_{\text{robust}}(T)] \ge \Gamma \cdot |\lambda_1 - \lambda_0| \cdot B_{\max} \cdot (T - \tau^*),
\label{eq:robust-regret}
\end{equation}
which is linear in the post-change horizon. The ratio $\mathbb{E}[R_{\text{adaptive}}(T)] / \mathbb{E}[R_{\text{robust}}(T)]$ approaches $\mathbb{E}[r_{\text{post}}] / (\Gamma \cdot |\lambda_1 - \lambda_0| \cdot B_{\max}) < 1$ as $T$ grows, establishing that the adaptive policy achieves strictly lower asymptotic regret.

\clearpage
\section{Simulation Architecture}
\label{sec:simulation}

We implement the ED simulation in SimPy, a process-based discrete-event simulation library for Python.

\subsection{ED Environment}

The \texttt{EDEnvironment} class manages all patient dynamics:
\begin{itemize}
\item \textbf{Bed resource:} A \texttt{simpy.Resource} with capacity $B_{\max} = 20$ beds.
\item \textbf{Arrival generation:} Inter-arrival times are drawn from $\operatorname{Exponential}(1/\lambda(t))$ where $\lambda(t)$ follows \eqref{eq:nhpp}. This produces a non-homogeneous Poisson process.
\item \textbf{Service times:} Log-normally distributed per acuity class.
\item \textbf{Queueing:} Separate FIFO queues per acuity level. On each arrival, the policy's \texttt{get\_action} method is called with the current state and determines admissions from the queue.
\item \textbf{Discharge:} When a patient completes service, the bed is released and the policy is re-invoked to admit from the queue.
\end{itemize}

\subsection{Scenarios}
\label{sec:scenarios}

We evaluate four scenarios designed to test different types of non-stationarity:

\begin{enumerate}
\item \textbf{Baseline} ($\tau^* = \infty$): No surge. Isolates the adaptive policy's false-alarm behavior under normal diurnal variation.
\item \textbf{Flash Flood (Step-Shift)}: At $t = 24$, $\lambda$ triples ($\times 3$) for 12 hours. Tests the CUSUM's ability to detect an abrupt, large-magnitude shift.
\item \textbf{Creeping Crisis (Linear Drift)}: At $t = 12$, $\lambda$ increases by a factor of 2.5 over 36 hours via hourly compounding at $1\%$ per hour. Tests detection of slowly evolving shifts where no single observation is extreme.
\item \textbf{Acuity Flip (Composition Shift)}: At $t = 24$, $\lambda$ doubles ($\times 2$) for 24 hours while the critical fraction rises from 10\% to 25\%. Tests joint detection of rate and acuity-composition changes.
\end{enumerate}

\subsection{Hindsight Oracle}
\label{sec:oracle}

The oracle provides a lower bound on the optimal total cost $V^*_T$, enabling regret computation. We implement two variants:

\textbf{Greedy heuristic:} Sort patients by acuity weight (descending) and arrival time (ascending), then assign each to the earliest available bed. This yields an upper bound on $V^*_T$.

\textbf{MILP formulation:} For small instances, we solve a time-indexed mixed-integer linear program:
\begin{align}
\min_{x_{i,t}} &\quad \sum_{i=1}^{N} w_i \sum_{t = \lceil a_i / R \rceil}^{T-1} (tR - a_i) x_{i,t} \\
\text{s.t.} &\quad \sum_{t} x_{i,t} = 1, \quad \forall i, \\
&\quad \sum_{i=1}^{N} \sum_{s = t - d_i + 1}^{t} x_{i,s} \le B_{\max}, \quad \forall t, \\
&\quad x_{i,t} \in \{0, 1\},
\end{align}
where $a_i$ is arrival time, $d_i = \lceil \text{service}_i / R\rceil$ is service duration in slots, $R$ is time resolution (default 1 hour), and $T = \lceil (\text{horizon} + \max_i \text{service}_i) / R\rceil$. We use HiGHS (open-source high-performance MILP solver) as the MILP solver with fallback to the greedy heuristic for large instances ($>20{,}000$ binary variables).

\section{Experimental Design}

\subsection{Parameter Calibration}

Arrival parameters are calibrated to synthetic NHAMCS data \cite{johnson2023mimic}. The NHPP in \eqref{eq:nhpp} uses $\lambda_{\text{base}} = 10.0$, $A = 0.3$, $\phi = \pi/2$, yielding a daily mean of 240 arrivals. Acuity composition: 10\% critical, 30\% urgent, 60\% routine, based on published ESI distributions. A MIMIC-IV demo cohort (236 ED visits) provides an alternative calibration with 42\% critical, 17\% urgent, 41\% routine, reflecting the higher-acuity profile of admitted ED patients.

\subsection{Simulation Parameters}

\begin{itemize}
\item Number of beds: $B_{\max} = 20$
\item Horizon: $T = 48$ hours (6 epochs of 8 hours each)
\item Robustness budget: $\Gamma = 0.2$
\item CUSUM threshold: $h = 5.0$ (calibrated for $\text{ARL}_0 \approx 168$ hours)
\item Sliding window: $w = 6$ observations
\item Replications: 30 per scenario with matched seeds across policies
\end{itemize}

\subsection{Statistical Methodology}

For each scenario and replication, we compute:
\begin{itemize}
\item Total cost $C^{\pi}$ for both robust and adaptive policies
\item Oracle cost $V^*$ (computed once per replication, shared across policies)
\item Regret $R^{\pi} = C^{\pi} - V^*$
\item Detection delay (adaptive only): time from surge onset to first CUSUM trigger
\end{itemize}

We report mean $\pm$ 95\% confidence interval across replications. Statistical significance is assessed via paired $t$-test and Wilcoxon signed-rank test, both comparing per-replication regret pairs.

\section{Results}

\subsection{Regret Comparison}

\begin{table*}[!htbp]
\centering
\caption{Cumulative regret comparison (mean $\pm$ 95\% CI, 30 replications).}
\label{tab:results}
\small
\begin{tabular}{lcccc}
\toprule
Scenario & Robust Regret & Adaptive Regret & Reduction & $p$-value \\
\midrule
Baseline      & $31\,013 \pm 6\,960$ & $21\,903 \pm 6\,099$ & 29.4\% & $<0.001$ \\
Flash Flood   & $53\,607 \pm 3\,882$ & $44\,205 \pm 3\,650$ & 17.5\% & $<0.001$ \\
Creeping Crisis & $46\,978 \pm 4\,481$ & $37\,738 \pm 4\,339$ & 19.7\% & $<0.001$ \\
Acuity Flip   & $42\,613 \pm 4\,471$ & $30\,694 \pm 4\,075$ & 28.0\% & $<0.001$ \\
\bottomrule
\end{tabular}
\end{table*}

Table~\ref{tab:results} presents the main results. The adaptive policy achieves statistically significant regret reduction across all four scenarios ($p < 0.001$ for all paired $t$-tests, confirmed by Wilcoxon signed-rank tests). The largest absolute regret occurs in the Flash Flood scenario (sudden rate spike), while the largest percentage improvement occurs in the Baseline and Acuity Flip scenarios.

A notable finding is that the adaptive policy improves over the robust policy even in the Baseline scenario (no engineered surge). This suggests the CUSUM detects and adapts to the natural diurnal rate variation, which ranges from approximately $\lambda_0(1 - A) = 7$ to $\lambda_0(1 + A) = 13$ across the daily cycle. The static robust policy, centered at $\lambda_0 = 10$, is systematically over-conservative during low-demand periods and under-protective during peaks.

\subsection{Detection Delay}

\begin{table}[!htbp]
\centering
\caption{Detection delay (mean hours). Lorden's bound is $1 + h / D_{\text{KL}}$ with $h = 5.0$.}
\label{tab:detect}
\small
\resizebox{\columnwidth}{!}{
\begin{tabular}{lccc}
\toprule
Scenario & Mean Delay (h) & Lorden Bound & Bound / Delay \\
\midrule
Flash Flood     & 1.13 & 3.17 & 2.80 \\
Creeping Crisis & 4.83 & 11.42 & 2.36 \\
Acuity Flip     & 2.47 & 5.89 & 2.38 \\
\bottomrule
\end{tabular}
}
\end{table}

Table~\ref{tab:detect} reports detection delays. The CUSUM detects the Flash Flood step shift rapidly (mean 1.13~hours), well within Lorden's theoretical bound of 3.17~hours. The bound-to-delay ratio of approximately 2.5 across all scenarios indicates the bound is conservative but not loose. The Creeping Crisis takes longer to detect (4.83~hours), consistent with the gradual drift and smaller per-step KL divergence.

Zero false alarms were recorded across all 30 baseline replications, confirming the threshold calibration ($h = 5.0$) is appropriately conservative.

\subsection{Regret Growth Over Time}

Figure~\ref{fig:cumulative_regret} shows cumulative regret trajectories (mean across 30 replications) for all four scenarios. In each case, the adaptive policy's regret grows more slowly than the robust policy's after the detection point, with the gap widening monotonically over the remaining horizon.

The widening gap is most pronounced in the Flash Flood and Acuity Flip scenarios, where the shift is large and abrupt, enabling rapid detection. In the Creeping Crisis scenario, the gap grows more slowly, consistent with the longer detection delay. Even in the Baseline scenario, the adaptive policy maintains a lower regret trajectory, suggesting the CUSUM corrects for diurnal variation that the static robust policy treats as noise.

The 48-hour horizon limits our ability to distinguish between logarithmic and asymptotic linear growth rates. A log-log regression of cumulative regret against time yields exponents in the range 0.6--0.8 for the adaptive policy, compared to 0.9--1.1 for the robust policy, consistent with sublinear growth for the adaptive policy. Longer experiments (168+ hours) would be needed to conclusively estimate the asymptotic growth exponent.

\begin{figure*}[!htbp]
\centering
\includegraphics[width=\textwidth]{fig_cumulative_regret.pdf}
\caption{Cumulative regret over time (mean across 30 replications). The adaptive policy (green) consistently diverges from the robust baseline (red) after the CUSUM detects the distribution shift. The gap widens over the remaining horizon, establishing that the adaptive policy achieves strictly lower cumulative regret.}
\label{fig:cumulative_regret}
\end{figure*}

\subsection{Statistical Significance}

\begin{table}[!htbp]
\centering
\caption{Paired statistical tests (adaptive vs.\ robust regret, 30 pairs).}
\label{tab:stats}
\small
\resizebox{\columnwidth}{!}{
\begin{tabular}{lccc}
\toprule
Scenario & $t$-statistic & $p$-value (paired $t$) & $p$-value (Wilcoxon) \\
\midrule
Baseline        & 8.40 & $<0.001$ & $<0.001$ \\
Flash Flood     & 11.98 & $<0.001$ & $<0.001$ \\
Creeping Crisis & 13.48 & $<0.001$ & $<0.001$ \\
Acuity Flip     & 24.22 & $<0.001$ & $<0.001$ \\
\bottomrule
\end{tabular}
}
\end{table}

All four scenarios reject the null hypothesis of equal means at $p < 0.001$ (Table~\ref{tab:stats}). The Wilcoxon signed-rank test confirms the results without normality assumptions. The Acuity Flip scenario produces the largest $t$-statistic (24.22), reflecting both the large percentage improvement and low variance.

\subsection{MIMIC-IV Validation}

To demonstrate generalizability beyond synthetic NHAMCS calibration, we repeat the Flash Flood experiment using parameters calibrated from the MIMIC-IV demo cohort \cite{johnson2023mimic}. The demo-calibrated acuity mix (42\% critical, 17\% urgent, 41\% routine) reflects a higher-acuity population than the NHAMCS baseline, since the demo dataset is sourced from patients subsequently admitted to the hospital.

\begin{table}[!htbp]
\centering
\caption{MIMIC-IV validation (Flash Flood, 48-hour horizon).}
\label{tab:mimic}
\small
\resizebox{\columnwidth}{!}{
\begin{tabular}{lcc}
\toprule
Metric & Robust & Adaptive \\
\midrule
Mean Regret & $38\,328$ & $30\,084$ \\
95\% CI & $\pm 6\,783$ & $\pm 6\,590$ \\
Reduction & \multicolumn{2}{c}{21.5\%} \\
Paired $t$ & \multicolumn{2}{c}{$t = 7.28$, $p < 0.001$} \\
\bottomrule
\end{tabular}
}
\end{table}

The adaptive policy maintains a statistically significant advantage (21.5\% reduction, $p < 0.001$) under real-data calibration (Table~\ref{tab:mimic}). The absolute regret values differ from the NHAMCS case due to the higher baseline critical acuity, which compresses the action space for both policies. The percentage reduction is comparable between the two calibrations, suggesting the framework's benefits are robust to the underlying acuity distribution.

\section{Discussion and Limitations}

\paragraph{Interpretation of results.} The empirical results support the core hypothesis: coupling CUSUM change detection with robust optimization reduces regret compared to a static robust policy. The mechanism is intuitive --- once the detector confirms a distributional shift, the policy re-centers its uncertainty set to reflect the new operating regime, rather than continuing to hedge against the original baseline. The magnitude of improvement (17--29\%) is practically meaningful in a healthcare context, where even small improvements in wait-time equity are clinically significant.

\paragraph{Oracle quality.} The hindsight oracle uses a greedy heuristic that is an upper bound on the true optimal cost, meaning our reported regret values are lower bounds on the true regret. Both policies are evaluated against the same oracle, so the \emph{relative} comparison remains valid. The MILP formulation using HiGHS provides a tighter bound for small instances but becomes computationally prohibitive for the full 30-replication experiment. Future work could employ decomposition methods or column generation to solve larger MILP instances.

\paragraph{Sublinearity.} The 48-hour horizon limits asymptotic analysis. While the empirical regret trajectories are qualitatively consistent with sublinear growth for the adaptive policy, longer experiments (168+ hours, as supported by \texttt{main.py}) would be needed to conclusively estimate the growth exponent. The theoretical bound from Lorden's inequality guarantees that detection delay is $\mathcal{O}(\log T)$ for thresholds $h = \Theta(\log T)$, but the post-detection regret scaling depends on the rate estimation accuracy and remains an open question for this specific setting.

\paragraph{Modeling assumptions.} Several modeling choices merit discussion. The action space excludes discharge decisions, which would affect bed availability dynamics in practice. The CUSUM operates on raw arrival counts rather than acuity-weighted counts, which limits sensitivity to composition shifts like the Acuity Flip scenario. The assumption of known post-shift magnitude $\lambda_1$ for CUSUM calibration is unrealistic in practice; adaptive CUSUM variants that estimate $\lambda_1$ online are a natural extension.

\section{Conclusion}

We presented an adaptive-regret framework for ED bed allocation that combines CUSUM change detection with online robust optimization. The framework detects distribution shifts via a Poisson log-likelihood ratio CUSUM and responds by re-centering and expanding its uncertainty set. Lorden's inequality provides a finite-sample bound on detection delay. In a SimPy ED simulation with NHAMCS-calibrated parameters, the adaptive policy achieves 17--29\% regret reduction over a static robust policy across four scenarios (all $p < 0.001$). MIMIC-IV validation confirms 21.5\% improvement under real-data-calibrated parameters. Detection delays are well within theoretical bounds, with rapid detection of step shifts (1.1 hours) and reasonable detection of gradual drifts (4.8 hours).

\paragraph{Future work.} Several extensions are natural. First, incorporating acuity-weighted CUSUM statistics would improve sensitivity to composition shifts. Second, adaptive CUSUM variants that estimate the post-shift magnitude online would relax the assumption of known $\lambda_1$. Third, extending the MDP formulation to include discharge decisions and stochastic service times would increase realism. Fourth, validation on the full MIMIC-IV-ED dataset (with ESI scores) would strengthen the empirical contribution. Fifth, the framework's modular design (change detector + robust optimizer) generalizes to other healthcare resource allocation problems, including operating room scheduling and ambulance fleet management.

\bibliographystyle{plainnat}
\bibliography{references}

\end{document}
